%% ODER: format ==         = "\mathrel{==}"
%% ODER: format /=         = "\neq "
%
%
\makeatletter
\@ifundefined{lhs2tex.lhs2tex.sty.read}%
  {\@namedef{lhs2tex.lhs2tex.sty.read}{}%
   \newcommand\SkipToFmtEnd{}%
   \newcommand\EndFmtInput{}%
   \long\def\SkipToFmtEnd#1\EndFmtInput{}%
  }\SkipToFmtEnd

\newcommand\ReadOnlyOnce[1]{\@ifundefined{#1}{\@namedef{#1}{}}\SkipToFmtEnd}
\usepackage{amstext}
\usepackage{amssymb}
\usepackage{stmaryrd}
\DeclareFontFamily{OT1}{cmtex}{}
\DeclareFontShape{OT1}{cmtex}{m}{n}
  {<5><6><7><8>cmtex8
   <9>cmtex9
   <10><10.95><12><14.4><17.28><20.74><24.88>cmtex10}{}
\DeclareFontShape{OT1}{cmtex}{m}{it}
  {<-> ssub * cmtt/m/it}{}
\newcommand{\texfamily}{\fontfamily{cmtex}\selectfont}
\DeclareFontShape{OT1}{cmtt}{bx}{n}
  {<5><6><7><8>cmtt8
   <9>cmbtt9
   <10><10.95><12><14.4><17.28><20.74><24.88>cmbtt10}{}
\DeclareFontShape{OT1}{cmtex}{bx}{n}
  {<-> ssub * cmtt/bx/n}{}
\newcommand{\tex}[1]{\text{\texfamily#1}}	% NEU

\newcommand{\Sp}{\hskip.33334em\relax}


\newcommand{\Conid}[1]{\mathit{#1}}
\newcommand{\Varid}[1]{\mathit{#1}}
\newcommand{\anonymous}{\kern0.06em \vbox{\hrule\@width.5em}}
\newcommand{\plus}{\mathbin{+\!\!\!+}}
\newcommand{\bind}{\mathbin{>\!\!\!>\mkern-6.7mu=}}
\newcommand{\rbind}{\mathbin{=\mkern-6.7mu<\!\!\!<}}% suggested by Neil Mitchell
\newcommand{\sequ}{\mathbin{>\!\!\!>}}
\renewcommand{\leq}{\leqslant}
\renewcommand{\geq}{\geqslant}
\usepackage{polytable}

%mathindent has to be defined
\@ifundefined{mathindent}%
  {\newdimen\mathindent\mathindent\leftmargini}%
  {}%

\def\resethooks{%
  \global\let\SaveRestoreHook\empty
  \global\let\ColumnHook\empty}
\newcommand*{\savecolumns}[1][default]%
  {\g@addto@macro\SaveRestoreHook{\savecolumns[#1]}}
\newcommand*{\restorecolumns}[1][default]%
  {\g@addto@macro\SaveRestoreHook{\restorecolumns[#1]}}
\newcommand*{\aligncolumn}[2]%
  {\g@addto@macro\ColumnHook{\column{#1}{#2}}}

\resethooks

\newcommand{\onelinecommentchars}{\quad-{}- }
\newcommand{\commentbeginchars}{\enskip\{-}
\newcommand{\commentendchars}{-\}\enskip}

\newcommand{\visiblecomments}{%
  \let\onelinecomment=\onelinecommentchars
  \let\commentbegin=\commentbeginchars
  \let\commentend=\commentendchars}

\newcommand{\invisiblecomments}{%
  \let\onelinecomment=\empty
  \let\commentbegin=\empty
  \let\commentend=\empty}

\visiblecomments

\newlength{\blanklineskip}
\setlength{\blanklineskip}{0.66084ex}

\newcommand{\hsindent}[1]{\quad}% default is fixed indentation
\let\hspre\empty
\let\hspost\empty
\newcommand{\NB}{\textbf{NB}}
\newcommand{\Todo}[1]{$\langle$\textbf{To do:}~#1$\rangle$}

\EndFmtInput
\makeatother
%
%
%
%
%
%
% This package provides two environments suitable to take the place
% of hscode, called "plainhscode" and "arrayhscode". 
%
% The plain environment surrounds each code block by vertical space,
% and it uses \abovedisplayskip and \belowdisplayskip to get spacing
% similar to formulas. Note that if these dimensions are changed,
% the spacing around displayed math formulas changes as well.
% All code is indented using \leftskip.
%
% Changed 19.08.2004 to reflect changes in colorcode. Should work with
% CodeGroup.sty.
%
\ReadOnlyOnce{polycode.fmt}%
\makeatletter

\newcommand{\hsnewpar}[1]%
  {{\parskip=0pt\parindent=0pt\par\vskip #1\noindent}}

% can be used, for instance, to redefine the code size, by setting the
% command to \small or something alike
\newcommand{\hscodestyle}{}

% The command \sethscode can be used to switch the code formatting
% behaviour by mapping the hscode environment in the subst directive
% to a new LaTeX environment.

\newcommand{\sethscode}[1]%
  {\expandafter\let\expandafter\hscode\csname #1\endcsname
   \expandafter\let\expandafter\endhscode\csname end#1\endcsname}

% "compatibility" mode restores the non-polycode.fmt layout.

\newenvironment{compathscode}%
  {\par\noindent
   \advance\leftskip\mathindent
   \hscodestyle
   \let\\=\@normalcr
   \let\hspre\(\let\hspost\)%
   \pboxed}%
  {\endpboxed\)%
   \par\noindent
   \ignorespacesafterend}

\newcommand{\compaths}{\sethscode{compathscode}}

% "plain" mode is the proposed default.
% It should now work with \centering.
% This required some changes. The old version
% is still available for reference as oldplainhscode.

\newenvironment{plainhscode}%
  {\hsnewpar\abovedisplayskip
   \advance\leftskip\mathindent
   \hscodestyle
   \let\hspre\(\let\hspost\)%
   \pboxed}%
  {\endpboxed%
   \hsnewpar\belowdisplayskip
   \ignorespacesafterend}

\newenvironment{oldplainhscode}%
  {\hsnewpar\abovedisplayskip
   \advance\leftskip\mathindent
   \hscodestyle
   \let\\=\@normalcr
   \(\pboxed}%
  {\endpboxed\)%
   \hsnewpar\belowdisplayskip
   \ignorespacesafterend}

% Here, we make plainhscode the default environment.

\newcommand{\plainhs}{\sethscode{plainhscode}}
\newcommand{\oldplainhs}{\sethscode{oldplainhscode}}
\plainhs

% The arrayhscode is like plain, but makes use of polytable's
% parray environment which disallows page breaks in code blocks.

\newenvironment{arrayhscode}%
  {\hsnewpar\abovedisplayskip
   \advance\leftskip\mathindent
   \hscodestyle
   \let\\=\@normalcr
   \(\parray}%
  {\endparray\)%
   \hsnewpar\belowdisplayskip
   \ignorespacesafterend}

\newcommand{\arrayhs}{\sethscode{arrayhscode}}

% The mathhscode environment also makes use of polytable's parray 
% environment. It is supposed to be used only inside math mode 
% (I used it to typeset the type rules in my thesis).

\newenvironment{mathhscode}%
  {\parray}{\endparray}

\newcommand{\mathhs}{\sethscode{mathhscode}}

% texths is similar to mathhs, but works in text mode.

\newenvironment{texthscode}%
  {\(\parray}{\endparray\)}

\newcommand{\texths}{\sethscode{texthscode}}

% The framed environment places code in a framed box.

\def\codeframewidth{\arrayrulewidth}
\RequirePackage{calc}

\newenvironment{framedhscode}%
  {\parskip=\abovedisplayskip\par\noindent
   \hscodestyle
   \arrayrulewidth=\codeframewidth
   \tabular{@{}|p{\linewidth-2\arraycolsep-2\arrayrulewidth-2pt}|@{}}%
   \hline\framedhslinecorrect\\{-1.5ex}%
   \let\endoflinesave=\\
   \let\\=\@normalcr
   \(\pboxed}%
  {\endpboxed\)%
   \framedhslinecorrect\endoflinesave{.5ex}\hline
   \endtabular
   \parskip=\belowdisplayskip\par\noindent
   \ignorespacesafterend}

\newcommand{\framedhslinecorrect}[2]%
  {#1[#2]}

\newcommand{\framedhs}{\sethscode{framedhscode}}

% The inlinehscode environment is an experimental environment
% that can be used to typeset displayed code inline.

\newenvironment{inlinehscode}%
  {\(\def\column##1##2{}%
   \let\>\undefined\let\<\undefined\let\\\undefined
   \newcommand\>[1][]{}\newcommand\<[1][]{}\newcommand\\[1][]{}%
   \def\fromto##1##2##3{##3}%
   \def\nextline{}}{\) }%

\newcommand{\inlinehs}{\sethscode{inlinehscode}}

% The joincode environment is a separate environment that
% can be used to surround and thereby connect multiple code
% blocks.

\newenvironment{joincode}%
  {\let\orighscode=\hscode
   \let\origendhscode=\endhscode
   \def\endhscode{\def\hscode{\endgroup\def\@currenvir{hscode}\\}\begingroup}
   %\let\SaveRestoreHook=\empty
   %\let\ColumnHook=\empty
   %\let\resethooks=\empty
   \orighscode\def\hscode{\endgroup\def\@currenvir{hscode}}}%
  {\origendhscode
   \global\let\hscode=\orighscode
   \global\let\endhscode=\origendhscode}%

\makeatother
\EndFmtInput
%



%% ----------- Applications --------------
%% %format applyx a b = a "@" b
%% %format applyx a b = a "\," b

%% Function application with no arguments

%% Function application with 1 argument

%% Function application with 1 argument, without parenthesis

%% Function application with many arguments

%% ----------- End of Applications --------------




%% Old version: | for BNF, [] for choice
%%%format choice = "[ \! ]"
%%%format `choice` = "\mathop{[ \! ]}"

%% New version: tall | for BNF, fat | for choice
%%format vbar = "\mathop{\bigm|}"
%%format choice = "\boldsymbol{|}"
%%format `choice` = "\mathop{\boldsymbol{|}}"




%% Arrays and tuples
%% %format array (x) = "\textbf{array}\{" x "\}"
%% %format arrKW = "\textbf{arr}"
%% %format arr (x) = arrKW "\{" x "\}"

%% %format defKW = "\textbf{def}"

% only used in comments, but still prettier this way:
%%%%format map = "\textbf{map}"
% format := = "\mapsto"
%  format effs x = "\!\!<\!\!" x "\!\!>\!\!" dots
























%% ----------- Metatheory --------------
%%%%%%%%%%%%%%%%%%%%%%%%%%%%%%%%%%%%

\newcommand{\rr}{\ensuremath{R}}
\newcommand{\go}[1]{{\rightarrow}_{#1}}
\newcommand{\gos}[1]{{\twoheadrightarrow}_{#1}}
\newcommand{\gok}[2]{{\twoheadrightarrow}_{#1}^{#2}}
\newcommand{\rstep}[3]{\ensuremath{#2 \go{#1} #3}}
\newcommand{\rsteps}[3]{\ensuremath{#2 \gos{#1} #3}}
\newcommand{\rstepk}[4]{\ensuremath{#3 \gok{#1}{#2} #4}}
\newcommand{\univ}{\ensuremath{\mathcal{E}}}
\newcommand{\join}[2]{{#1} \sim {#2}}
\newcommand{\dia}{\ensuremath{\diamond}}

\section{Preliminaries} \label{sec:preliminaries}

\begin{definition}[Reductions and Normal Forms] A \emph{reduction} \rr is a binary relation on a set of terms \univ.
We write $\rr^{k}$ for the $k$-step closure of $\rr$ and $\rr^{*}$ for the reflexive
and transitive closure of $\rr$, \ie $\rr^{*} \equiv \cup_{0 \leq k} \rr^k$.
%
We write
%
$\rstep{\rr}{a}{b}$ ($a$ \emph{steps to} $b$) if $(a, b) \in \rr$,
%
$\rstepk{\rr}{k}{a}{b}$ ($a$ \emph{$k$-steps to} $b$) if $(a, b) \in \rr^{k}$, and
%
$\rsteps{\rr}{a}{b}$ ($a$ \emph{reduces to} $b$) if $(x, y) \in \rr^{*}$.
%
%
A term $a$ is an $\rr$-\emph{Normal Form} if there does not exist any $b$
such that $\rstep{\rr}{a}{b}$.
\end{definition}

For clarity, we will omit the reduction $\rr$ when it is clear from the context.

\begin{definition}[Diamond Property and Confluence]
%
A reduction satisfies the \emph{diamond} property if whenever
$\rstep{}{a}{b}$ and $\rstep{}{a}{c}$, there is a $d$ such
that $\rstep{}{b}{d}$ and $\rstep{}{c}{d}$.
%
We say that two terms $b$, $c$ can be \emph{joined}, written $\join{b}{c}$,
if there is a $d$ such that $\rsteps{}{b}{d}$ and $\rsteps{}{c}{d}$.
%
A reduction is \emph{locally confluent} if whenever
$\rstep{}{a}{b}$ and $\rstep{}{a}{c}$, we have $\join{b}{c}$.
%
A reduction is \emph{confluent} if whenever
$\rsteps{}{a}{b}$ and $\rsteps{}{a}{c}$, we have $\join{b}{c}$.
\end{definition}

\begin{figure}[h]
\begin{minipage}[t]{0.30\textwidth}
% https://q.uiver.app/?q=WzAsNCxbMSwwLCJ4Il0sWzAsMSwieV8xIl0sWzIsMSwieV8yIl0sWzEsMiwieiJdLFswLDJdLFswLDFdLFsxLDMsIiIsMix7InN0eWxlIjp7ImJvZHkiOnsibmFtZSI6ImRvdHRlZCJ9fX1dLFsyLDMsIiIsMix7InN0eWxlIjp7ImJvZHkiOnsibmFtZSI6ImRvdHRlZCJ9fX1dXQ==
\[\begin{tikzcd}
	& a \\
	b && {c} \\
	& d
	\arrow[from=1-2, to=2-3]
	\arrow[from=1-2, to=2-1]
	\arrow[dotted, from=2-1, to=3-2]
	\arrow[dotted, from=2-3, to=3-2]
\end{tikzcd}\]
\end{minipage}
\begin{minipage}[t]{0.30\textwidth}
% https://q.uiver.app/?q=WzAsNCxbMSwwLCJ4Il0sWzAsMSwieV8xIl0sWzIsMSwieV8yIl0sWzEsMiwieiJdLFswLDJdLFswLDFdLFsxLDMsIiIsMix7InN0eWxlIjp7ImJvZHkiOnsibmFtZSI6ImRvdHRlZCJ9LCJoZWFkIjp7Im5hbWUiOiJlcGkifX19XSxbMiwzLCIiLDIseyJzdHlsZSI6eyJib2R5Ijp7Im5hbWUiOiJkb3R0ZWQifSwiaGVhZCI6eyJuYW1lIjoiZXBpIn19fV1d
\[\begin{tikzcd}
	& a \\
	b && c \\
	& d
	\arrow[from=1-2, to=2-3]
	\arrow[from=1-2, to=2-1]
	\arrow[dotted, two heads, from=2-1, to=3-2]
	\arrow[dotted, two heads, from=2-3, to=3-2]
\end{tikzcd}\]
\end{minipage}
\begin{minipage}[t]{0.30\textwidth}
% https://q.uiver.app/?q=WzAsNCxbMSwwLCJ4Il0sWzAsMSwieV8xIl0sWzIsMSwieV8yIl0sWzEsMiwieiJdLFswLDIsIiIsMCx7InN0eWxlIjp7ImhlYWQiOnsibmFtZSI6ImVwaSJ9fX1dLFswLDEsIiIsMCx7InN0eWxlIjp7ImhlYWQiOnsibmFtZSI6ImVwaSJ9fX1dLFsxLDMsIiIsMix7InN0eWxlIjp7ImJvZHkiOnsibmFtZSI6ImRvdHRlZCJ9LCJoZWFkIjp7Im5hbWUiOiJlcGkifX19XSxbMiwzLCIiLDIseyJzdHlsZSI6eyJib2R5Ijp7Im5hbWUiOiJkb3R0ZWQifSwiaGVhZCI6eyJuYW1lIjoiZXBpIn19fV1d
\[\begin{tikzcd}
	& a \\
	b && c \\
	& d
	\arrow[two heads, from=1-2, to=2-3]
	\arrow[two heads, from=1-2, to=2-1]
	\arrow[dotted, two heads, from=2-1, to=3-2]
	\arrow[dotted, two heads, from=2-3, to=3-2]
\end{tikzcd}\]
\end{minipage}
\caption{\textbf{Diamond Property (L), Local Confluence (M), and Confluence (R)}}
\end{figure}

% [Barendregt 1976]
\begin{lemma}[Diamond-Confluence] \label{lem:diamond}
    If $\rr$ satisfies the diamond property then $\rr$ is confluent.
\end{lemma}

% [Barendregt 1976] [Huet 1980]
\begin{lemma}[Confluent-NF] \label{lem:nf}
    If $\rr$ is confluent then every term reduces to at most one normal form.
\end{lemma}

% [Barendregt 1976] [Huet 1980]
\begin{lemma}[Confluent-Closure] \label{lem:closure}
    If $\rr$ is confluent then $\rr^{*}$ is confluent.
\end{lemma}




\mypara{Commutativity}
%
Two reductions $R$ and $S$ are \emph{commutative} if for all terms
$a, b, c$ such that $\rsteps{R}{a}{b}$ and $\rsteps{S}{a}{c}$
there exists $d$ such that $\rsteps{S}{b}{d}$ and $\rsteps{R}{c}{d}$,
as illustrated on the left in \cref{fig:commutativity}.
%
Two reductions $R$ and $S$ are \emph{strongly commutative} if for all terms
$a, b, c$ such that $\rstep{R}{a}{b}$ and $\rstep{R}{a}{c}$
there exists $d$ such that $\rstep{S}{b}{d}$ and $\rstep{R}{c}{d}$,
as illustrated on the right in \cref{fig:commutativity}.
%$R \cup S$ satisfies the diamond property.

\begin{figure}[h]
%
\begin{minipage}[t]{0.45\textwidth}
%% COMMUTATIVE
% https://q.uiver.app/?q=WzAsNCxbMSwwLCJhIl0sWzAsMSwiYiJdLFsyLDEsImMiXSxbMSwyLCJkIl0sWzAsMiwiUyIsMCx7InN0eWxlIjp7ImhlYWQiOnsibmFtZSI6ImVwaSJ9fX1dLFswLDEsIlIiLDIseyJzdHlsZSI6eyJoZWFkIjp7Im5hbWUiOiJlcGkifX19XSxbMSwzLCJSIiwyLHsic3R5bGUiOnsiYm9keSI6eyJuYW1lIjoiZG90dGVkIn0sImhlYWQiOnsibmFtZSI6ImVwaSJ9fX1dLFsyLDMsIlMiLDAseyJzdHlsZSI6eyJib2R5Ijp7Im5hbWUiOiJkb3R0ZWQifSwiaGVhZCI6eyJuYW1lIjoiZXBpIn19fV1d
\[\begin{tikzcd}
	& a \\
	b && c \\
	& d
	\arrow["S", two heads, from=1-2, to=2-3]
	\arrow["R"', two heads, from=1-2, to=2-1]
	\arrow["S"', dotted, two heads, from=2-1, to=3-2]
	\arrow["R", dotted, two heads, from=2-3, to=3-2]
\end{tikzcd}\]
\end{minipage}
\begin{minipage}[t]{0.45\textwidth}
% https://q.uiver.app/?q=WzAsNCxbMSwwLCJhIl0sWzAsMSwiYiJdLFsyLDEsImMiXSxbMSwyLCJkIl0sWzAsMiwiUyJdLFswLDEsIlIiLDJdLFsxLDMsIlIiLDIseyJzdHlsZSI6eyJib2R5Ijp7Im5hbWUiOiJkb3R0ZWQifX19XSxbMiwzLCJTIiwwLHsic3R5bGUiOnsiYm9keSI6eyJuYW1lIjoiZG90dGVkIn19fV1d
\[\begin{tikzcd}
	& a \\
	b && c \\
	& d
	\arrow["S", from=1-2, to=2-3]
	\arrow["R"', from=1-2, to=2-1]
	\arrow["S"', dotted, from=2-1, to=3-2]
	\arrow["R", dotted, from=2-3, to=3-2]
\end{tikzcd}\]
\end{minipage}
\caption{\textbf{Commutativity (L) and Strong Commutativity (R)}}
\label{fig:commutativity}
\end{figure}

% [Huet 1980]
\begin{lemma}[Strong-Commutativity] If $R$ and $S$ strongly commute then $R$ and $S$ commute.
\end{lemma}

\begin{proof}
Via the following ``chase'' diagram.

    % https://q.uiver.app/?q=WzAsMTYsWzAsMCwiXFxidWxsZXQiXSxbMCwxLCJcXGJ1bGxldCJdLFswLDIsIlxcYnVsbGV0Il0sWzAsMywiXFxidWxsZXQiXSxbMSwwLCJcXGJ1bGxldCJdLFsyLDAsIlxcYnVsbGV0Il0sWzMsMCwiXFxidWxsZXQiXSxbMSwxLCJcXGJ1bGxldCJdLFsyLDEsIlxcYnVsbGV0Il0sWzMsMSwiXFxidWxsZXQiXSxbMSwyLCJcXGJ1bGxldCJdLFsyLDIsIlxcYnVsbGV0Il0sWzMsMiwiXFxidWxsZXQiXSxbMSwzLCJcXGJ1bGxldCJdLFsyLDMsIlxcYnVsbGV0Il0sWzMsMywiXFxidWxsZXQiXSxbMCw0LCJTIl0sWzQsNSwiUyIsMCx7InN0eWxlIjp7ImJvZHkiOnsibmFtZSI6ImRhc2hlZCJ9fX1dLFs1LDYsIlMiXSxbMCwxLCJSIiwyXSxbMSwyLCJSIiwyLHsic3R5bGUiOnsiYm9keSI6eyJuYW1lIjoiZGFzaGVkIn19fV0sWzIsMywiUiIsMl0sWzEsNywiUyIsMCx7InN0eWxlIjp7ImJvZHkiOnsibmFtZSI6ImRvdHRlZCJ9fX1dLFs0LDcsIlIiLDIseyJzdHlsZSI6eyJib2R5Ijp7Im5hbWUiOiJkb3R0ZWQifX19XSxbNyw4LCJTIiwwLHsic3R5bGUiOnsiYm9keSI6eyJuYW1lIjoiZGFzaGVkIn19fV0sWzgsOSwiUyIsMCx7InN0eWxlIjp7ImJvZHkiOnsibmFtZSI6ImRvdHRlZCJ9fX1dLFs1LDgsIlIiLDIseyJzdHlsZSI6eyJib2R5Ijp7Im5hbWUiOiJkb3R0ZWQifX19XSxbNiw5LCJSIiwyLHsic3R5bGUiOnsiYm9keSI6eyJuYW1lIjoiZG90dGVkIn19fV0sWzksMTIsIlIiLDIseyJzdHlsZSI6eyJib2R5Ijp7Im5hbWUiOiJkYXNoZWQifX19XSxbOCwxMSwiUiIsMix7InN0eWxlIjp7ImJvZHkiOnsibmFtZSI6ImRhc2hlZCJ9fX1dLFs3LDEwLCJSIiwyLHsic3R5bGUiOnsiYm9keSI6eyJuYW1lIjoiZGFzaGVkIn19fV0sWzIsMTAsIlMiLDAseyJzdHlsZSI6eyJib2R5Ijp7Im5hbWUiOiJkb3R0ZWQifX19XSxbMywxMywiUyIsMCx7InN0eWxlIjp7ImJvZHkiOnsibmFtZSI6ImRvdHRlZCJ9fX1dLFsxMywxNCwiUyIsMCx7InN0eWxlIjp7ImJvZHkiOnsibmFtZSI6ImRhc2hlZCJ9fX1dLFsxMCwxMywiUiIsMix7InN0eWxlIjp7ImJvZHkiOnsibmFtZSI6ImRvdHRlZCJ9fX1dLFsxMSwxNCwiUiIsMix7InN0eWxlIjp7ImJvZHkiOnsibmFtZSI6ImRvdHRlZCJ9fX1dLFsxMiwxNSwiUiIsMix7InN0eWxlIjp7ImJvZHkiOnsibmFtZSI6ImRvdHRlZCJ9fX1dLFsxNCwxNSwiUyIsMCx7InN0eWxlIjp7ImJvZHkiOnsibmFtZSI6ImRvdHRlZCJ9fX1dLFsxMCwxMSwiUyIsMCx7InN0eWxlIjp7ImJvZHkiOnsibmFtZSI6ImRhc2hlZCJ9fX1dLFsxMSwxMiwiUyIsMCx7InN0eWxlIjp7ImJvZHkiOnsibmFtZSI6ImRvdHRlZCJ9fX1dXQ==
\[\begin{tikzcd}
	\bullet & \bullet & \bullet & \bullet \\
	\bullet & \bullet & \bullet & \bullet \\
	\bullet & \bullet & \bullet & \bullet \\
	\bullet & \bullet & \bullet & \bullet
	\arrow["S", from=1-1, to=1-2]
	\arrow["S", dashed, from=1-2, to=1-3]
	\arrow["S", from=1-3, to=1-4]
	\arrow["R"', from=1-1, to=2-1]
	\arrow["R"', dashed, from=2-1, to=3-1]
	\arrow["R"', from=3-1, to=4-1]
	\arrow["S", dotted, from=2-1, to=2-2]
	\arrow["R"', dotted, from=1-2, to=2-2]
	\arrow["S", dashed, from=2-2, to=2-3]
	\arrow["S", dotted, from=2-3, to=2-4]
	\arrow["R"', dotted, from=1-3, to=2-3]
	\arrow["R"', dotted, from=1-4, to=2-4]
	\arrow["R"', dashed, from=2-4, to=3-4]
	\arrow["R"', dashed, from=2-3, to=3-3]
	\arrow["R"', dashed, from=2-2, to=3-2]
	\arrow["S", dotted, from=3-1, to=3-2]
	\arrow["S", dotted, from=4-1, to=4-2]
	\arrow["S", dashed, from=4-2, to=4-3]
	\arrow["R"', dotted, from=3-2, to=4-2]
	\arrow["R"', dotted, from=3-3, to=4-3]
	\arrow["R"', dotted, from=3-4, to=4-4]
	\arrow["S", dotted, from=4-3, to=4-4]
	\arrow["S", dashed, from=3-2, to=3-3]
	\arrow["S", dotted, from=3-3, to=3-4]
\end{tikzcd}\]
\end{proof}

% Classical "commutativity lemma"
\begin{lemma}[Commutativity] \label{lem:comm}
If $R$ and $S$ are confluent and commute, then $R \cup S$ is confluent.
\end{lemma}

\begin{lemma}[Union] \label{lem:comm-union}
    If $R$ and $S_1$ commute and $R$ and $S_2$ commute then $R$ and $S_1 \cup S_2$ commute.
\end{lemma}

\begin{proof}
Via the following chase diagram (probably a standard result.)

% https://q.uiver.app/?q=WzAsMTIsWzAsMCwiXFxidWxsZXQiXSxbMCwxLCJcXGJ1bGxldCJdLFsxLDAsIlxcYnVsbGV0Il0sWzIsMCwiXFxidWxsZXQiXSxbMywwLCJcXGJ1bGxldCJdLFsxLDEsIlxcYnVsbGV0Il0sWzIsMSwiXFxidWxsZXQiXSxbMywxLCJcXGJ1bGxldCJdLFs0LDAsIlxcYnVsbGV0Il0sWzUsMCwiXFxidWxsZXQiXSxbNCwxLCJcXGJ1bGxldCJdLFs1LDEsIlxcYnVsbGV0Il0sWzAsMiwiU18xIiwwLHsic3R5bGUiOnsiaGVhZCI6eyJuYW1lIjoiZXBpIn19fV0sWzIsMywiU18yIiwwLHsic3R5bGUiOnsiaGVhZCI6eyJuYW1lIjoiZXBpIn19fV0sWzMsNCwiIiwwLHsic3R5bGUiOnsiYm9keSI6eyJuYW1lIjoiZGFzaGVkIn0sImhlYWQiOnsibmFtZSI6ImVwaSJ9fX1dLFswLDEsIlIiLDIseyJzdHlsZSI6eyJoZWFkIjp7Im5hbWUiOiJlcGkifX19XSxbMSw1LCJTXzEiLDAseyJzdHlsZSI6eyJib2R5Ijp7Im5hbWUiOiJkb3R0ZWQifSwiaGVhZCI6eyJuYW1lIjoiZXBpIn19fV0sWzIsNSwiUiIsMix7InN0eWxlIjp7ImJvZHkiOnsibmFtZSI6ImRvdHRlZCJ9LCJoZWFkIjp7Im5hbWUiOiJlcGkifX19XSxbNSw2LCJTXzIiLDAseyJzdHlsZSI6eyJib2R5Ijp7Im5hbWUiOiJkb3R0ZWQifSwiaGVhZCI6eyJuYW1lIjoiZXBpIn19fV0sWzYsNywiIiwwLHsic3R5bGUiOnsiYm9keSI6eyJuYW1lIjoiZG90dGVkIn0sImhlYWQiOnsibmFtZSI6ImVwaSJ9fX1dLFszLDYsIlIiLDIseyJzdHlsZSI6eyJib2R5Ijp7Im5hbWUiOiJkb3R0ZWQifSwiaGVhZCI6eyJuYW1lIjoiZXBpIn19fV0sWzQsNywiUiIsMix7InN0eWxlIjp7ImJvZHkiOnsibmFtZSI6ImRvdHRlZCJ9LCJoZWFkIjp7Im5hbWUiOiJlcGkifX19XSxbNCw4LCJTXzEiLDAseyJzdHlsZSI6eyJoZWFkIjp7Im5hbWUiOiJlcGkifX19XSxbOCw5LCJTXzIiLDAseyJzdHlsZSI6eyJoZWFkIjp7Im5hbWUiOiJlcGkifX19XSxbOCwxMCwiUiIsMix7InN0eWxlIjp7ImJvZHkiOnsibmFtZSI6ImRvdHRlZCJ9LCJoZWFkIjp7Im5hbWUiOiJlcGkifX19XSxbOSwxMSwiUiIsMix7InN0eWxlIjp7ImJvZHkiOnsibmFtZSI6ImRvdHRlZCJ9LCJoZWFkIjp7Im5hbWUiOiJlcGkifX19XSxbNywxMCwiU18xIiwwLHsic3R5bGUiOnsiYm9keSI6eyJuYW1lIjoiZG90dGVkIn0sImhlYWQiOnsibmFtZSI6ImVwaSJ9fX1dLFsxMCwxMSwiU18yIiwwLHsic3R5bGUiOnsiYm9keSI6eyJuYW1lIjoiZG90dGVkIn0sImhlYWQiOnsibmFtZSI6ImVwaSJ9fX1dXQ==
\[\begin{tikzcd}
	\bullet & \bullet & \bullet & \bullet & \bullet & \bullet \\
	\bullet & \bullet & \bullet & \bullet & \bullet & \bullet
	\arrow["{S_1}", two heads, from=1-1, to=1-2]
	\arrow["{S_2}", two heads, from=1-2, to=1-3]
	\arrow[dashed, two heads, from=1-3, to=1-4]
	\arrow["R"', two heads, from=1-1, to=2-1]
	\arrow["{S_1}", dotted, two heads, from=2-1, to=2-2]
	\arrow["R"', dotted, two heads, from=1-2, to=2-2]
	\arrow["{S_2}", dotted, two heads, from=2-2, to=2-3]
	\arrow[dotted, two heads, from=2-3, to=2-4]
	\arrow["R"', dotted, two heads, from=1-3, to=2-3]
	\arrow["R"', dotted, two heads, from=1-4, to=2-4]
	\arrow["{S_1}", two heads, from=1-4, to=1-5]
	\arrow["{S_2}", two heads, from=1-5, to=1-6]
	\arrow["R"', dotted, two heads, from=1-5, to=2-5]
	\arrow["R"', dotted, two heads, from=1-6, to=2-6]
	\arrow["{S_1}", dotted, two heads, from=2-4, to=2-5]
	\arrow["{S_2}", dotted, two heads, from=2-5, to=2-6]
\end{tikzcd}\]
\end{proof}

\begin{lemma}[N-Commutativity] \label{lem:comm-n}
If
%
(i)~$\forall 0 \leq i \leq n$, the reduction $R_i$ is confluent, and
%
(ii)~$\forall 0 \leq i < j \leq n$, the reductions $R_i$ and $R_j$ commute
%
then $\cup_{i=0}^n R_i$ is confluent.
\end{lemma}

\begin{proof}
    By induction on $n$ using \cref{lem:comm} and \cref{lem:comm-union}.
\end{proof}
